\section{Abundance after nonvanishing}
\label{bl:sec:abundance}

Over $\C$, the supported theorem closes the zero-Iitaka-dimension
case once lower-dimensional abundance generates the whole reduced
floor. For positive Iitaka dimension, the induction also needs
effective minimal models of fiber adjoints, which need not be nef.
We organize these two uses in one proposition, then verify its
hypotheses in dimension four and under the all-dimensional premise.

Every effective representative below is effective up to actual
$\Q$-linear equivalence. Ordinary log minimal models supply nef
adjoints; semiampleness is the conclusion to be proved. The first
lemma records exactly which section spaces survive passage to such
a model.

\subsection{Minimal models and the whole reduced floor}

\begin{lemma}[Comparison with an ordinary log minimal model]
\label{bl:lem:abundance-model-comparison}
Let $(W,\Delta_W)$ be a projective lc rational pair and let
$(V,\Delta_V)$ be a $\Q$-factorial dlt log minimal model.  Use the
log birational model convention in which divisors extracted on $V$
have coefficient one.  On a common resolution
$\alpha:U\to W$, $\beta:U\to V$, with compatible canonical divisors,
\begin{equation}
 \alpha^*(K_W+\Delta_W)
   =\beta^*(K_V+\Delta_V)+E,
 \qquad E\geq0,\quad E\text{ is }\beta\text{-exceptional}.
 \label{bl:eq:abundance-model-comparison}
\end{equation}
Consequently the adjoints have the same section spaces in sufficiently
divisible degrees, and hence the same Kodaira dimension.
\end{lemma}

\begin{proof}
Set $E$ equal to the difference in
\eqref{bl:eq:abundance-model-comparison}.  It is anti-nef over $W$, since
$K_V+\Delta_V$ is nef.  The discrepancy condition for a log minimal
model gives $\alpha_*E\geq0$: only prime divisors of $W$ contracted
by the model map can contribute, and their discrepancies improve.
The negativity lemma therefore gives $E\geq0$.

At a prime divisor on $V$ which is also a divisor on $W$, the
coefficient of $E$ is zero.  At a prime divisor on $V$ extracted over
$W$, it is the new log discrepancy minus the old log discrepancy.
The former is zero and the latter is nonnegative.  Its coefficient
is thus at most zero, and the effectivity just proved makes it zero.
This proves exceptionality over $V$.

For divisible $m$, adding the effective exceptional divisor $mE$
to $\beta^*m(K_V+\Delta_V)$ creates no additional sections.
Indeed a rational function allowed by the former divisor has the
required pole bounds at every prime divisor on the normal variety
$V$, and therefore is a section downstairs; the converse follows
by pullback.  Birational pullback from $W$ preserves sections as
well.  These identifications prove the assertion.
\end{proof}

We may take the $\Q$-factorial dlt models in this lemma in the
ordinary minimal-model existence statements below.  A crepant dlt
modification, when needed, has the same adjoint pullback and preserves
nefness and all section spaces.  The dlt modification theorem is,
for example, \cite[Theorem~2.8]{FujinoGongyo2014}.

\begin{lemma}[Semiampleness on the whole floor]
\label{bl:lem:abundance-floor}
Assume abundance for projective lc rational pairs of dimension less
than $d$.  If $(V,C)$ is a projective $\Q$-factorial dlt pair of
dimension $d$ and $A=K_V+C$ is nef, then the actual restricted
rational line bundle $A|_P$, where $P=\lfloor C\rfloor$, is
semiample.  Its restriction to any reduced union of components of
$P$ is also semiample.
\end{lemma}

\begin{proof}
There is nothing to prove if $P$ is empty.  Whole-floor dlt
adjunction gives an effective rational boundary $\Delta_P$ for
which $(P,\Delta_P)$ is semi-divisorial log terminal, in particular
semi-log-canonical, and
\begin{equation}
 K_P+\Delta_P=A|_P.
 \label{bl:eq:abundance-floor-adjunction}
\end{equation}
Here the equality is an adjunction identification of rational line
bundles on the entire reduced scheme $P$; see
\cite[Remark~2.7]{FujinoGongyo2014}.  In particular, it specifies the
gluing along the double locus.

For clarity, if $\zeta:\coprod P_i\to P$ is the normalization,
its pullback is the disjoint union of the adjunction formulas
$K_{P_i}+\Delta_i=A|_{P_i}$.  The conductor has coefficient one in
$\Delta_i$ and the normalization pairs are dlt.  The whole-floor
statement includes the $S_2$ and nodal-in-codimension-one properties
of $P$.  The adjunction maps are pluriresidue maps: at a generic
double point, the two branch residues have the usual opposite
subsequent residues of an ambient logarithmic form.  Taking an
even divisible power removes the sign convention and gives the
dualizing-sheaf gluing.  Thus \eqref{bl:eq:abundance-floor-adjunction}
is not merely a collection of componentwise equivalences; it
identifies a divisible restricted invertible sheaf on $P$ itself.

Each $K_{P_i}+\Delta_i$ is nef and is semiample by the assumed
lower-dimensional abundance.  Semiampleness descends from the
normalization of a projective slc pair by
\cite[Theorem~4.3, equivalently Theorem~1.5]{FujinoGongyo2014}.
Applied to \eqref{bl:eq:abundance-floor-adjunction}, this gives
generation of a power on the whole $P$.  Restricting its evaluation
map to a reduced closed union of components preserves surjectivity,
which proves the last assertion.  This also treats extra
coefficient-one components meeting the selected union and its strata.
\end{proof}

\subsection{The induction step}

\begin{proposition}[Abundance from effective models and lower-dimensional abundance]
\label{bl:prop:abundance-induction}
Fix $d\geq1$ and assume the following two statements over $\C$:
\begin{enumerate}
 \item every projective lc rational pair of dimension at most $d$
 whose adjoint has an effective $\Q$-linearly equivalent divisor
 has an ordinary log minimal model;
 \item every nef projective lc rational adjoint in dimension less
 than $d$ is semiample.
\end{enumerate}
Then, for a projective lc rational pair $(X,B_X)$ of dimension $d$,
\[
 D=K_X+B_X\text{ nef},\qquad \kappa(X,D)\geq0
 \quad\Longrightarrow\quad D\text{ semiample}.
\]
If $\kappa(X,D)=0$, then $D\sim_{\Q}0$.
\end{proposition}

\begin{proof}
We first treat Kodaira dimension zero, then the Iitaka fibers in
the positive-dimensional case, and finally the lc centers.  The
second assumption is full lower-dimensional abundance; no
nonvanishing hypothesis is imposed on the lower-dimensional
restrictions in Lemma~\ref{bl:lem:abundance-floor}.

\smallskip\noindent
\emph{Kodaira dimension zero.}
Choose $M\geq0$ with $D\sim_{\Q}M$.  Suppose that $M\ne0$.
Take a projective log resolution $p:W\to X$ of the pair together
with $\Supp M$.  Start with the strict transform of $B_X$, raise
the coefficient to one along every strict transform of a component
of $M$, and put coefficient one on every exceptional prime divisor.
Call the resulting snc boundary $C_W$, and let $T_W\geq0$ be the
sum of the coefficient increases on nonexceptional divisors.
Using log discrepancies, we have
\begin{equation}
 \begin{split}
 K_W+C_W&\sim_{\Q}M_W,\\
 M_W&=p^*M+\sum_{E_i\text{ exceptional}}
       a(E_i;X,B_X)E_i+T_W\geq0.
 \end{split}
 \label{bl:eq:abundance-raised-boundary}
\end{equation}
The pair $(W,C_W)$ is lc, and
$\Supp M_W\subset\Supp\lfloor C_W\rfloor$.

Raising the coefficients here preserves Kodaira dimension zero.
To see this directly, a rational function in any divisible system
$|mM_W|$ has, on $X$, poles only along $\Supp M$.  There is a
constant $c>0$, independent of $m$, such that the permitted pole
order at every such prime is at most $cm$ times its coefficient
in $M$: this is a finite list of positive coefficients.  Increasing
$c$ and clearing denominators identifies this rational function
with a section of some $|nM|$.  Since $M$ is effective and
$\kappa(M)=0$, every such section is constant in the rational-function
description of $H^0(X,\OO_X(nM))$.  Thus
$h^0(W,mM_W)=1$ for every sufficiently divisible $m$, and
$\kappa(W,K_W+C_W)=0$.

By the first assumption, take an ordinary log minimal model
$(V_*,C_*)$ of $(W,C_W)$.  On a common resolution
$\alpha:U\to W$, $\beta:U\to V_*$ set
\[
 M_*:=\beta_*\alpha^*M_W.
\]
Lemma~\ref{bl:lem:abundance-model-comparison} shows that
$M_*\geq0$, $M_*\sim_{\Q}K_{V_*}+C_*$, and that the latter
adjoint has Kodaira dimension zero.  The support of $M_*$ lies
in $\lfloor C_*\rfloor$: surviving components of $M_W$ already
have coefficient one, and any component extracted over $W$ has
coefficient one by the model convention.

Crucially, $M_*\ne0$.  Otherwise $\alpha^*M_W$ would be
$\beta$-exceptional.  The nonzero effective divisor
\[
 Q=(p\alpha)^*M\leq\alpha^*M_W
\]
would then be exceptional over $V_*$ as well.  But
$Q\sim_{\Q}(p\alpha)^*D$ is nef.  The negativity lemma applied
to this effective exceptional, relatively nef divisor forces
$Q=0$, a contradiction.  This is the place where nefness of the
original adjoint ensures that the effective representative
survives on the new model.

Choose $q>0$ sufficiently divisible that $G=qM_*$ is Cartier and
$G\sim q(K_{V_*}+C_*)$.  Lemma~\ref{bl:lem:abundance-floor} makes
$\OO_{V_*}(G)|_{G_{\mathrm{red}}}$ semiample.  All the supported
pair hypotheses now hold, with $G\ne0$.  Proposition~\ref{bl:prop:complex-supported}
gives $\kappa(V_*,K_{V_*}+C_*)>0$, a contradiction.  Therefore
$M=0$ and $D\sim_{\Q}0$.

\smallskip\noindent
\emph{A non-nef adjoint on an Iitaka fiber.}
Put $k=\kappa(X,D)>0$.  Resolve an Iitaka map of a sufficiently
divisible multiple of $D$ and take its Stein factorization,
using a simultaneous projective log resolution $\pi:W\to X$:
\[
 f:W\longrightarrow Z,\qquad \dim Z=k.
\]
With $\Gamma$ the strict transform of $B_X$ plus the reduced
exceptional divisor, there is an actual equality
\begin{equation}
 L:=K_W+\Gamma=\pi^*D+N,
 \qquad N\geq0\text{ exceptional over }X.
 \label{bl:eq:abundance-nonnef-adjoint}
\end{equation}
In particular $\kappa(W,L)=k$.  The resolved linear system also
gives, for a positive integer $u$ and an ample rational divisor
$A_Z$,
\begin{equation}
 uL\sim_{\Q}f^*A_Z+E_0,\qquad E_0\geq0.
 \label{bl:eq:abundance-iitaka-bound}
\end{equation}
The ample divisor is obtained by pulling back the hyperplane
bundle through the finite map in the Stein factorization.

Let $F$ be a very general smooth fiber.  Generic smoothness for
the finitely many horizontal snc strata and avoidance of the
vertical strata give the log smooth lc pair $(F,\Gamma_F)$, with
\[
 L|_F=K_F+\Gamma_F.
\]
It has an effective rational representative: restrict any fixed
nonzero divisible section of $L$, choosing $F$ outside the locus
where the section vanishes identically.  Moreover
$\kappa(F,L|_F)=0$.  Here is a justification that does not require
$L|_F$ to be nef.  Work with a Cartier multiple $\mathcal L$ of
$L$ and clear \eqref{bl:eq:abundance-iitaka-bound} to write
$v\mathcal L\sim f^*A+E_1$, with $A$ ample Cartier and $E_1\geq0$.
Choose $F$ outside the countable union of the proper closed
sets needed for base change of $f_*\OO_W(t\mathcal L)$, for
all $t>0$.  If some fiber system has positive-dimensional image,
then, after twisting that direct image by a sufficiently positive
multiple $aA$, its global sections generate the fiber system at
the general point of $Z$.  Multiplying by sections of a further
very ample multiple $bA$ gives a system in
\[
 |t\mathcal L+(a+b)f^*A|
\]
which detects both the base and a nonconstant fiber ratio.  Its
image has dimension at least $k+1$.  Multiplication by the section
of $(a+b)E_1$ embeds this system in
$|(t+(a+b)v)\mathcal L|$, contradicting $\kappa(W,L)=k$.

Assume $s:=\dim F>0$.  The pair $(F,\Gamma_F)$ has dimension
$s=d-k<d$ and an effective adjoint, so it has an ordinary log
minimal model $(F_*,\Gamma_*)$ by the first assumption.
Its nef adjoint has Kodaira dimension zero by
Lemma~\ref{bl:lem:abundance-model-comparison}.  The second assumption
therefore makes it semiample, and hence $\Q$-linearly trivial.
Notice that this applies lower-dimensional abundance on $F_*$;
the divisor $K_F+\Gamma_F$ on the original fiber may be non-nef.

On a smooth common resolution $\tau:\widetilde F\to F$ and
$\mu:\widetilde F\to F_*$, the comparison gives
\begin{equation}
 \begin{gathered}
 P+N_F\sim_{\Q}E_F,\qquad P:=\tau^*(\pi^*D|_F)\text{ nef},\\
 N_F:=\tau^*(N|_F)\geq0,\qquad
 E_F\geq0\text{ exceptional over }F_*.
 \end{gathered}
 \label{bl:eq:abundance-fiber-comparison}
\end{equation}
If $H$ is the pullback of an ample divisor on $F_*$, the
projection formula and effectivity yield the full inequality
\begin{equation}
 0\leq P\cdot H^{s-1}
       =-N_F\cdot H^{s-1}\leq0.
 \label{bl:eq:abundance-fiber-sandwich}
\end{equation}
Since $H$ is nef and big, write $H\sim_{\Q}\epsilon A+T$ with
$A$ ample, $\epsilon>0$, and $T\geq0$.  Intersections of $T$
with products of nef classes are nonnegative.  Replacing one
factor $H$ at a time in the zero intersection in
\eqref{bl:eq:abundance-fiber-sandwich} thus gives
$P\cdot A^{s-1}=0$.  The Hodge index theorem for a nef class
then gives $P\equiv0$; when $s=1$ this is the degree criterion.
It follows that $\pi^*D|_F\equiv0$.  For zero-dimensional $F$
the same conclusion is automatic.

\smallskip\noindent
\emph{Numerical dimension and lc centers.}
For a nef divisor, $\kappa(D)\leq\nu(D)$.  We prove the reverse
inequality here.  If $\nu(D)>k$, a general sufficiently ample
complete intersection $T\subset W$ of dimension $k+1$ has
\[
 (\pi^*D|_T)^{k+1}>0.
\]
The restriction is nef and big.  The map $T\to Z$ is dominant
and its very general fiber is an integral curve, by Bertini
on the generic fiber of $W\to Z$.  The degree of $\pi^*D|_T$
on this curve is zero by the preceding argument.  On the other
hand, a big decomposition
$\pi^*D|_T\sim_{\Q}\epsilon A_T+E_T$, with $A_T$ ample and
$E_T\geq0$, gives strictly positive degree on a very general
fiber curve: the curves cover $T$, so the general one is not
contained in $\Supp E_T$.  This contradiction proves
$\nu(X,D)=\kappa(X,D)$.

Take a projective crepant $\Q$-factorial dlt modification
$\xi:(X',B')\to(X,B_X)$.  Its adjoint $D'=\xi^*D$ is nef
and abundant.  Lemma~\ref{bl:lem:abundance-floor} makes its
restriction to $\lfloor B'\rfloor$ semiample.  Every lc center
of the dlt pair is contained in this floor; restriction to the
center and then pullback to its normalization remain semiample.
Thus $D'$ is log abundant in the precise sense of
\cite[Definition~4.1]{FujinoGongyo2014}.  The nef and log abundant
semiampleness theorem
\cite[Theorem~4.2, equivalently Theorem~1.6]{FujinoGongyo2014}
applies to $(X',B')$.  Finally, semiampleness descends along
$\xi$, since $\xi_*\OO_{X'}=\OO_X$ and a divisible Cartier
multiple of $D'$ is the pullback of that of $D$.  This proves
the proposition.
\end{proof}

\subsection{The two applications}

\begin{proof}[Proof of Theorem~\ref{thm:abundance-after-nonvanishing} over $\C$]
Projective lc abundance through dimension three is the classical
threefold theorem \cite[Theorem~1.1]{KMM1994}, with its correction
\cite{KMM2004}.  A direct surface reference is
\cite[Theorem~7.2]{Fujino2012}; curves are elementary.  These
statements concern effective rational boundaries and the actual
log canonical divisor.

Ordinary log minimal models exist for effective lc fourfolds by
\cite[Corollary~1.6]{Birkar2010}.  Equivalently, one may use
\cite[Corollary~1.7]{Birkar2011} together with the known
\emph{relative} three-dimensional minimal-model theorem for
pseudo-effective $\Q$-factorial dlt pairs.  The relative theorem,
not merely absolute threefold abundance, is the premise of that
corollary.  Its effectivity hypothesis is satisfied by the
$\Q$-linear effectivity used here.  The lower-dimensional
minimal-model statements are known as well.

Proposition~\ref{bl:prop:abundance-induction} with $d=4$ now proves
semiampleness under $\kappa\geq0$.  In Kodaira dimension zero
its proof gives $D\sim_{\Q}0$.  The lower-dimensional cases
follow from the cited abundance theorems.  The complex assertion is now complete. Section~\ref{bl:sec:fields}
proves its field extension as a separate corollary.
\end{proof}

\begin{proof}[Proof of Theorem~\ref{bl:thm:conditional} over $\C$]
Assume that every smooth projective complex variety $W$ with
pseudo-effective $K_W$ has a nonzero section of $mK_W$ for some
$m>0$.  In particular $K_W$ is not required to be nef.
Hashizume's theorem \cite[Theorem~1.4]{Hashizume2018} yields
lc nonvanishing and existence of ordinary log minimal models
in all dimensions.  The smooth premise is exactly Conjecture~1.3
of that paper; the output is its Conjectures~1.1 and~1.2.

For rational data, its real-linear nonvanishing gives rational-linear
nonvanishing. Apply the independent finite rational-feasibility lemma
of \cite[Lemma~8.1]{OpenAINonvanishing2026}
to the finite coefficients of $D$ and of the principal divisors in an
expression $D\sim_{\R}M_{\R}\geq0$. That lemma uses only rational
linear algebra, and no nonvanishing or abundance theorem. It produces
$M_{\Q}\geq0$ with $D\sim_{\Q}M_{\Q}$ on the same normal variety;
denominators may be cleared together with any specified Cartier index.

Induct on the dimension, starting with dimension zero.  For a
nef lc adjoint in dimension $d$, Hashizume supplies nonvanishing,
because nefness implies pseudo-effectivity, and supplies all
the effective ordinary minimal models required through dimension
$d$.  The induction hypothesis supplies full nef lc abundance
in every smaller dimension.  Proposition~\ref{bl:prop:abundance-induction}
therefore applies.  This proves the conditional abundance
statement over $\C$, without adding semiampleness to Hashizume's
minimal-model conclusion.  The ground-field transfer follows next.
\end{proof}
